% Part: first-order-logic
% Chapter: sequent-calculus
% Section: quantifier-rules

\documentclass[../../../include/open-logic-section]{subfiles}

\begin{document}

\olfileid[tr]{fol}{seq}{qrl}

\olsection{Niceleyici Kuralları}

\subsection{$\lforall$ için Kurallar}

\begin{defish}
\Axiom$ !A(t), \Gamma \fCenter \Delta$
\RightLabel{\LeftR{\lforall}}
\UnaryInf$ \lforall[x][!A(x)],\Gamma \fCenter \Delta$
\DisplayProof
\hfill
\Axiom$ \Gamma \fCenter \Delta, !A(a) $
\RightLabel{\RightR{\lforall}}
\UnaryInf$ \Gamma \fCenter \Delta, \lforall[x][!A(x)]$
\DisplayProof
\end{defish}

\LeftR{\lforall} kuralında $t$ kapalı bir terimdir (yani değişken içermez).
\RightR{\lforall} kuralında $a$, \RightR{\lforall} kuralının alt sekantında
hiçbir yerde bulunmaması gereken !!a{constant}{}tir. $a$'ya
\RightR{\lforall} çıkarımının \emph{özdeğişkeni} deriz.\footnote{Yukarıdaki
kuralda $a$ aslında !!a{constant} olmasına rağmen ``özdeğişken'' terimini
kullanıyoruz. Bunun tarihsel nedenleri vardır.}

\subsection{$\lexists$ için Kurallar}

\begin{defish}
\Axiom$ !A(a), \Gamma \fCenter \Delta $
\RightLabel{\LeftR{\lexists}}
\UnaryInf$ \lexists[x][!A(x)], \Gamma \fCenter \Delta$
\DisplayProof
\hfill
\Axiom$ \Gamma \fCenter \Delta, !A(t) $
\RightLabel{\RightR{\lexists}}
\UnaryInf$ \Gamma \fCenter \Delta, \lexists[x][!A(x)]$
\DisplayProof
\end{defish}

Burada da $t$ kapalı bir terim, $a$ ise \LeftR{\lexists} kuralının alt
sekantında bulunmayan !!a{constant}{}tir. $a$'ya \LeftR{\lexists} çıkarımının
\emph{özdeğişkeni} deriz.

Bir özdeğişkenin \RightR{\lforall} veya \LeftR{\lexists} çıkarımının alt
sekantında bulunmaması koşuluna \emph{özdeğişken koşulu} denir.

\begin{explain}
$!A$ bir !!a{formula}, $x$ de bu formülde serbest bir !!{variable} olsun. Bu
durumu $!A(x)$ yazarak belirttiğimiz uzlaşımı hatırlayın. Aynı bağlamda $!A(t)$,
$\Subst{!A}{t}{x}$ ifadesinin kısaltmasıdır. Dolayısıyla
\RightR{\lexists} kuralını şöyle de yazabiliriz:
\begin{prooftree}
  \Axiom$\Gamma \fCenter \Delta, \Subst{!A}{t}{x}$
  \RightLabel{\RightR{\lexists}}
  \UnaryInf$\Gamma \fCenter \Delta, \lexists[x][!A]$
\end{prooftree}
$t$'nin $!A$ içinde zaten bulunabileceğine dikkat edin; örneğin $!A$,
$\Atom{\Obj P}{t,x}$ olabilir. Bu nedenle $\Gamma \Sequent \Delta,
\Atom{\Obj P}{t,t}$ sekantından $\Gamma \Sequent \Delta,
\lexists[x][\Atom{\Obj P}{t,x}]$ sekantını çıkarmak
\RightR{\lexists} kuralının doğru bir uygulamasıdır: öncüldeki $t$
oluşumlarından birini veya birkaçını, hepsini değiştirmek zorunda olmadan,
bağlı !!{variable}~$x$ ile ``değiştirebilirsiniz.'' Ancak
\RightR{\lforall} ile \LeftR{\lexists} kurallarındaki özdeğişken koşulları,
!!{constant}~$a$'nın $!A$ içinde bulunmamasını gerektirir. Dolayısıyla
$\Gamma \Sequent \Delta, \Atom{\Obj P}{a,a}$ sekantından
$\Gamma \Sequent \Delta, \lforall[x][\Atom{\Obj P}{a,x}]$ sekantını
$\RightR{\lforall}$ kullanarak doğru biçimde çıkaramazsınız.
\end{explain}

\begin{explain}
\RightR{\lexists} ile \LeftR{\lforall} kurallarında $t$ terimine herhangi bir
kısıtlama getirilmez. Buna karşılık \LeftR{\lexists} ile
\RightR{\lforall} kurallarında özdeğişken koşulu, !!{constant}~$a$'nın üst
sekantta $!A(a)$ içindeki $x$ yerine geçen oluşumları dışında hiçbir yerde
bulunmamasını gerektirir. Bu koşul sistemin sağlamlığı için gereklidir: bu
sayede yalnızca geçerli sekantları !!{derive}s mümkündür. Bu
koşul olmasaydı aşağıdaki çıkarımlara izin verilirdi:
\begin{prooftree}
  \Axiom$!A(a) \fCenter !A(a)$
  \RightLabel{*\LeftR{\lexists}}
  \UnaryInf$\lexists[x][!A(x)] \fCenter !A(a)$
  \RightLabel{\RightR{\lforall}}
  \UnaryInf$\lexists[x][!A(x)] \fCenter \lforall[x][!A(x)]$
  \DisplayProof\bottomAlignProof
  \qquad
  \Axiom$!A(a) \fCenter !A(a)$
  \RightLabel{*\RightR{\lforall}}
  \UnaryInf$!A(a) \fCenter \lforall[x][!A(x)]$
  \RightLabel{\LeftR{\lexists}}
  \UnaryInf$\lexists[x][!A(x)] \fCenter \lforall[x][!A(x)]$
\end{prooftree}
Oysa $\lexists[x][!A(x)] \Sequent \lforall[x][!A(x)]$ geçerli değildir.
\end{explain}

\end{document}
